% !TeX root = ../tesis.tex
%%%%%%%%%%%%%%%%%%%%%%%%%%%%%%%%%%%%%%%%%%%%%%%%%%%%%%%%%%%%%%%%%%%
% SECCIÓN 7.2 — EJE URBANO (Issue #23)
% Papel en el hilo: lo que la medición de la estructura revela.
% Esquema: diagnóstico → lectura social → acercar no es conectar →
% radial a anillo-radial y modo → CETRAM y lectura histórica →
% aporte, pendientes y transición a 7.3. Plan: CONTEXTO_CAP7.md §5.
% Las cifras de la pregunta de investigación quedan aquí; la entrada
% solo da el veredicto y remite a esta sección (decisión 2026-11-03).
% Fuera a propósito: Moscú/Melbourne, $408 M (1958), 92.7 % de
% superficie; «dos ejes, red menos frágil» solo como pendiente.
%%%%%%%%%%%%%%%%%%%%%%%%%%%%%%%%%%%%%%%%%%%%%%%%%%%%%%%%%%%%%%%%%%%
\section{Eje Urbano}
\label{sec:conclusiones-urbano}

Medida la estructura, la red de la \zmvm{} muestra un patrón que la captación
por corredor no registraba: eficiencia interna y exclusión perimetral
(§\ref{subsubsec:espectro-cierre}). La red produce viajes eficientes para quien
ya forma parte de ella; el problema está en quién queda fuera. La cobertura es
de 13.29\,\% de la \zmvm{}, con 53.52\,\% en la \cdmx{} y 3.92\,\% en los
municipios conurbados, 42 de los cuales no tienen ninguna estación a distancia
caminable (§\ref{subsec:cobertura-espacial}). La mitad de los nodos tiene siete
conexiones o menos. El tiempo de viaje promedio es de 113.91~min, y la
intermediación se concentra en Tacubaya ($B(v) = 0.2232$) y en un corredor que
va de Garibaldi a Barranca del Muerto, con un Gini de 0.8007
(§\ref{subsec:fase3-centralidad}). La red, en suma, no tiene un problema de
operación, sino de distribución.

Esa asimetría tiene una lectura social. La \zmvm{} se puede ordenar en tres
contornos de acceso al transporte; en el tercero, que es la corona periférica,
llegar al Metro exige dos o tres transbordos en Corredores Concesionados,
\textsc{rtp} o servicios no formalizados, de modo que el viaje cuesta tiempo y
dinero a la vez. \citet{rendon2009ciudadanos} habla de ciudadanos de primera,
segunda y tercera clase. La clase la define el transporte más cercano. Para
esos habitantes el traslado es también una forma de \termino{pobreza de tiempo}
\citep{gaceta_cuidados_2026}, porque el viaje se come el margen que queda para
el descanso y el cuidado. Sobre la \zmvm{}, entonces, no cabe un juicio único:
una sola red resulta eficiente para una parte del territorio y deficiente para
otra.

El escenario del anillo se lee sobre esa línea base. Su primera lección es que
acercar no es conectar (§\ref{subsec:anillo-conclusiones-contornos}). La
cobertura apenas pasa de 13.29 a 13.43\,\% en los dos casos y los municipios
sin estación cercana bajan de 42 a 41, porque el Periférico corre por la parte
de la ciudad que ya tiene más servicio. Lo que sí se mueve es la conexión de
quien ya está cerca de la red: el tiempo de viaje promedio baja 4.1\,\% con el
Anillo MB y 10.1\,\% con el Anillo Metro (§\ref{subsec:anillo-tiempo-viaje}).
Nezahualcóyotl, Chimalhuacán y Ecatepec ganan entre media hora y una hora con
cualquiera de los dos. El sur depende del modo. Tlalpan, Xochimilco, Tláhuac y
La Magdalena Contreras ganan seis o siete minutos con el primero y entre 15 y
27 con el segundo.

La segunda lección es que el anillo cambia la forma de la red, pero no reparte
su carga (§\ref{subsec:anillo-conclusiones-centralidad}). Con el Anillo MB,
Tacubaya pierde 3.1\,\% de su intermediación y los diez nodos principales siguen
en el centro. Con el Anillo Metro pierde 29.7\,\%, y nueve de esos diez lugares
pasan a estaciones del Periférico. La red deja de ser solo radial y se vuelve
anillo-radial, con un segundo eje sobre el arco sur y oriente. El Gini, sin
embargo, sube a 0.8134: la concentración cambia de lugar, no se reparte. La
diferencia entre los casos la marca el modo
(§\ref{subsec:anillo-conclusiones-modo}). A unos 14~km/h, el Anillo MB ocupa el
papel de la \textsc{rtp} que hoy circula por el Periférico; a 36~km/h, el Anillo
Metro compite con las líneas radiales y reordena la jerarquía. Por eso el
Caso~3 es el más favorable.

El Caso~3 deja además un tipo de nodo que la red actual no tiene. Seis
estaciones del anillo, entre ellas Constitución de 1917, Canal Nacional y
Barranca del Muerto, alcanzan una intermediación de 0.12 a 0.16 con solo 6 a 12
conexiones: por ellas pasarían muchos caminos, pero llegan pocas rutas
(§\ref{subsec:anillo-conclusiones-cetram}). La recomendación es planearlas como
\textsc{cetram} de nueva generación, con rutas alimentadoras y accesos de
última milla. Hay también una lectura histórica. El Periférico se construyó para
el automóvil en los años en que el Metro se rechazaba por costoso
(§\ref{subsec:sesgo-vial} y §\ref{subsec:rechazo-metro}), y en 1985 sus dos
corredores salieron del Programa Maestro por captación. Los resultados sugieren
que esa vialidad es la que mejor podría alojar un transporte masivo lateral, si
tiene vía exclusiva. Eso era lo que no se veía cuando solo se medía la demanda.

El aporte urbano de la tesis es, primero, un diagnóstico de la red
metropolitana completa, con la \cdmx{} y el Estado de México juntos, que queda
como línea base pública y reproducible. Segundo, precisa lo que suele decirse
en general sobre la periferia: lo que falta ahí es cobertura y capilaridad más
que velocidad, y la dependencia de la red tiene nombres y lugares. Tercero,
muestra que el modo y el derecho de vía cambian el tipo de efecto de un
corredor, no solo su tamaño. Quedan preguntas abiertas. El modelo no dice
cuántas personas viven cerca de cada estación, cuánto costaría cada opción, qué
condiciones físicas tiene el trazo ni si dos ejes harían a la red menos frágil
(§\ref{subsec:corredor-anillar-limitaciones}). Tampoco resuelve un problema de
escala. Con el crecimiento disperso de la \zmvm{}, el Periférico ya no es la
periferia: hoy marca el borde exterior del Primer y el Segundo Contorno
(Figura~\ref{fig:semovi-contornos}), y la periferia actual pediría un anillo
exterior. Lo que el modelo no mide define la agenda de
§\ref{sec:conclusiones-academico}.
